\section{Обоснование выбора IAPWS-IF97 как основы модели}

В работе принят стандарт IAPWS-IF97 (редакция 2012 года) как нормативная
и физическая основа вычислительной модели~\cite{iapws_if97_2012}.
Выбор обусловлен промышленной направленностью стандарта,
формализованной топологией области расчёта
и наличием обратных зависимостей для практических маршрутов.
По сравнению с табличными подходами и библиотеками общего назначения
IF97 даёт нормативно определённый,
достаточно быстрый
и инженерно интерпретируемый базис для построения собственного вычислительного ядра.

IF97 задаёт уравнения для регионов 1--5, линию насыщения (регион~4)
и правила перехода между областями, включая границу B23.
Это позволяет построить алгоритм, в котором регион определяется автоматически,
а вычисления выполняются по профильной модели выбранной области.

Стандарт не содержит единой универсальной аналитической формулы для всех входных пар
во всей рабочей области.
Для региона~3 используются специализированные обратные зависимости,
в том числе релиз для $v(p,T)$~\cite{iapws_vpt_region3}.
Маршрут $(\rho, T)$ реализуется как численное расширение.
Таким образом, IF97 остаётся неизменной физической базой,
а численные методы используются как инженерное дополнение,
не нарушающее нормативную интерпретацию результата.

\section{Требования к программному комплексу}

Требования формулируются исходя из инженерных сценариев и ограничений IF97.
Ключевой принцип: единый источник вычислений и одинаковый результат
во всех каналах доставки.

Функциональные требования:
\begin{itemize}
\item расчёт состояния воды и водяного пара в области применимости IAPWS-IF97;
\item поддержка маршрутов $pT$, $ph$, $ps$, $px$, $\rho T$ по входным параметрам $p$, $T$, $h$, $s$, $\rho$, $x$;
\item автоматическое определение региона IF97 с учётом линии насыщения и границы B23;
\item возврат унифицированного набора выходных свойств ($v$, $\rho$, $h$, $s$, $u$, $c_p$, $w$) и номера региона;
\item валидация диапазонов и корректная диагностика ошибок ввода и расчёта;
\item единый контракт обмена через \texttt{if97\_app\_api} для пользовательских приложений и сервисного контура.
\end{itemize}

Нефункциональные требования:
\begin{itemize}
\item вычислительная эффективность для интерактивного режима и пакетных расчётов;
\item воспроизводимость сборки и детерминированность результата;
\item устойчивость на фазовых границах и в околокритической области с явной сигнализацией о несходимости;
\item отсутствие дублирования предметной математики в пользовательских и сервисных слоях;
\item возможность масштабирования пакетной обработки и контроля нагрузки;
\item расширяемость архитектуры за счёт разделения ядра и тонких адаптерных слоёв;
\item сохранение физической интерпретации IF97 при численных расширениях (в частности, для $\rho T$).
\end{itemize}

Архитектура реализуется по схеме «одно вычислительное ядро и несколько каналов доставки».
Ядро \texttt{if97\_core} выполняет расчёт и валидацию.
Слой \texttt{if97\_app\_api} фиксирует DTO и контракты обмена.
Поверх ядра работают FLTK-клиент,
клиент на базе Yew, WebAssembly и Tauri
и gRPC-сервис,
что исключает расхождение формул между компонентами.

В качестве эксплуатационного ориентира для сервисного канала принимается
пропускная способность порядка $10^5$ точек/с
в пакетных gRPC-сценариях.
Актуальный контрольный прогон
подтверждает данный порядок величины:
среднее значение составило 129~632,83~точек/с.
Этот ориентир используется как проверяемый критерий при развитии системы.

\section{Постановка вычислительной задачи и поддерживаемые режимы ввода-вывода}

Вычислительная задача формулируется как отображение
из допустимой входной пары параметров в полное состояние рабочей среды.
Множество входных параметров: $p$, $T$, $h$, $s$, $\rho$, $x$.
Множество выходных параметров: $v$, $\rho$, $h$, $s$, $u$, $c_p$, $w$, номер региона IF97.

Для каждой допустимой входной пары требуется:
\begin{itemize}
\item определить регион IF97 и корректный маршрут расчёта;
\item вычислить унифицированный набор выходных свойств;
\item вернуть диагностируемую ошибку при недопустимом вводе или выходе за границы применимости.
\end{itemize}

Отдельное требование относится к региону~4.
Система должна корректно обрабатывать линию насыщения и двухфазные состояния,
поскольку эти режимы напрямую связаны с практическими задачами кипения,
конденсации и оценки влажности пара.

Поддерживаемые режимы ввода-вывода (маршруты расчёта) включают:
\begin{itemize}
\item прямой маршрут $pT$: расчёт состояния по давлению и температуре;
\item обратные маршруты $ph$ и $ps$: расчёт по давлению и энтальпии или энтропии;
\item маршрут $px$: расчёт по давлению и степени сухости в двухфазной области;
\item маршрут $\rho T$: численное расширение на основе метода секущих
с подбором давления и последующим вычислением по $pT$.
\end{itemize}

Во всех режимах обязательны валидация входа и контроль области применимости.
При выходе за границы IF97 или при некорректной комбинации параметров
возвращается диагностируемая ошибка, а не неопределённые значения.

Критерием корректности постановки является
совпадение физической модели и формата результата
для всех каналов доставки
(\texttt{if97\_core}, FLTK,
связки Yew, WebAssembly и Tauri,
gRPC).
Критерием эксплуатационной применимости является
сохранение целевого порядка производительности
в пакетных сценариях и в контейнерном запуске.

Сформулированная постановка задачи определяет логику следующих глав:
математическую модель и численные методы (глава 2),
реализацию \texttt{if97\_core} (глава 3),
клиентские приложения (глава 4)
и сервисно-инфраструктурный контур (глава 5).
